% !PRINTED 634
\weberpage{634}
We now pass to the substitution $t$, and in
\weq{2}{(\varepsilon, x) = x_0 + \sum_{1,\,n-1}^{h}\varepsilon^h x_h}
put $x_{gh}$ for $x_h$; this function then goes over into
\weq{3}{x_0 + \sum_{1,\,n}^{h}\varepsilon^h x_{gh}.} % [sic] Print has lower limits $1,n$ although the residue system excludes $0$; preserved.
In this sum $h$ may run through an arbitrary system of residues with respect to the modulus $n$, with the exclusion of $0$, and we can therefore replace $h$ by $g^{-1}h$; thereby the sum (3) becomes
\[
x_0 + \sum^{h}\varepsilon^{g^{-1}h}x_h = (\varepsilon^{g^{-1}}, x).
\]
If one applies the permutations $s$ and $t$ repeatedly, there results the influence of $s^\lambda, t^\lambda$, which consists in the permutations
\[
(\varepsilon, x), \varepsilon^{-\lambda}(\varepsilon, x)
\]
\[
(\varepsilon, x), (\varepsilon^{g^{-\lambda}}, x)
\]
and so in particular, for $\lambda = -1$,
\[
(\varepsilon, x), \varepsilon(\varepsilon, x)
\]
\[
(\varepsilon, x), (\varepsilon^{g}, x),
\]
and we obtain the following theorem:

\gesperrt{4. The generating permutations of the metacyclic group, $s, t$, effect the permutations
\[
(s)\colon (\varepsilon, x), \varepsilon^{-1}(\varepsilon, x); \qquad (s^{-1})\colon (\varepsilon, x), \varepsilon(\varepsilon, x)
\]
\[
(t)\colon (\varepsilon, x), (\varepsilon^{g^{-1}}, x); \qquad (t^{-1})\colon (\varepsilon, x), (\varepsilon^{g}, x).
\]}

But from 2. and 4., because $s$ is the generating permutation of the cyclic group, there results the consequence:

\gesperrt{5. If one represents the functions
\weq{4}{(\varepsilon, x)^n = \Phi(\varepsilon), \quad (\varepsilon^\lambda, x)(\varepsilon, x)^{-\lambda} = F(\varepsilon)}
in the normal form [\S.\,183, (3)], then the coefficients are cyclic functions of $x_0, x_1 \ldots x_{n-1}$ (\S.\,163).}

Here $\lambda$ may denote any arbitrary integer that is not divisible by $n$.

From the function $F(\varepsilon)$ there arise, if we take $\lambda = g$ and put for $\varepsilon$ its $n-1$ different values, $n-1$ functions, which all remain unchanged under $s$, and which we shall denote as follows:
%CONT>>
